% LSPV-3 REPORT::TOTALENV1::2026-10-03T02:38:00-07:00
\documentclass[11pt]{article}
\usepackage[margin=.88in]{geometry}
\usepackage{amsmath,amssymb}
\begin{document}
\section*{LSPV-3 TOTALENV1: full polynomial-exponent reconstruction envelope}
\textbf{STATUS: OPEN.} P versus NP remains open. This report proves deterministic accounting identities only. It proves no SAT lower bound, no arbitrary-algorithm acquisition theorem, and no barrier-crossing hinge.
\section*{CORPUS REVIEWED}
The current repository corpus was surveyed before substantive work. Governance read includes AGENTS.md, CODEX.md, and the pnp-research skill. Root main.tex and verified MAIN 0058 were read. Current task surfaces, materially relevant MAIN history, LSPV-1 structural work through FULLPAIR6/PAIRDEG1, LSPV-2 EXCESS/EXCESSDUAL1/TURANEX1, LSPV-3 FACTOR1/QLEN1/RAMCOMP1/ENVELOPE1/RAMDOM1/CROSSDOM1/CERTREG1, and LSPV-4 TRACECLOSE1/LEGALFLAT1/RESPSEL1 were reconciled. HEC, VSC-HEC, G-PPC-HEC, CRS-PPC, CCSPC and current hostile controls were surveyed. The immediate LSPV-2 predecessor surface remains verified; no exact preserved failed peer scientific body requiring recovery was found.
Primary literature was refreshed. Viola's alternating-time simulations and oracle obstruction remain relativizing. Van Melkebeek--Mocelin Sdroievski remains the targeted-HSG/AM white-box source. Akshay et al. prove Skolem/Herbrand extraction from proofs in a declared QBF proof system; proof acquisition is not free here.
\section*{DUPLICATION / CONTRADICTION AUDIT}
ENVELOPE1 gives three alternative route exponents
\[D(K)=\max\{a,q+bK\},\quad R(K)=\ell+c_0K^2,\quad F(K)=\max\{f,c_FK\},\]
\[E(K)=\min\{D(K),R(K),F(K)\}.\]
RAMDOM1 and CROSSDOM1 give pairwise dominance regions. CERTREG1 proves monotonicity under looser certificates for the same fixed decider. The prior ledger writes common charges additively but does not state the exact polynomial-exponent reduction of the entire charged sum. TOTALENV1 supplies that reduction. No contradiction was found.
\section*{TOTALENV1 exact theorem}
Let every mandatory common reconstruction charge satisfy
\[\tau_j(N)\le N^{g_j+o(1)}\]
for the fixed finite set of Q acquisition, normalization, census/output, amplification, indexing, common compilation, output writing, and transport charges. Let
\[G=\max_j g_j.\]
Since the number of common charge classes is constant,
\[\sum_j N^{g_j+o(1)}=N^{G+o(1)}.\]
Executing the cheapest available correct route after the common work gives
\[\boxed{T_{\rm total}(N)\le N^{\max\{G,E(K)\}+o(1)}}.\]
Equivalently,
\[\boxed{T_{\rm total}(N)\le N^{\max\{G,\min[\max\{a,q+bK\},\ell+c_0K^2,\max\{f,c_FK\}]\}+o(1)}}.\]
This is \textbf{PROVED/EXACT} as polynomial-exponent accounting under the declared route and common-charge hypotheses.
A route improvement is invisible at total exponent scale whenever \(G\ge E(K)\). Thus reducing a route from a quadratic-in-K exponent to a linear-in-K exponent does not improve the total polynomial exponent unless the route remains the active bottleneck.
Common charges cannot be moved inside the route minimum merely to improve the exponent. A cost is common only if every declared implementation must pay the same SAME\_REALIZATION obligation. Route-specific evaluator acquisition, compilation, or indexing remains inside the corresponding branch.
\section*{Quantifier and reconstruction audit}
The load-bearing predicate remains
\[\exists\sigma\;\forall u\;\exists z\;\forall y\;R(x,Q,\sigma,u,z,y),\]
with one shared size-normalized Q, Arthur randomness, and exact SAT YES/NO conditions typed in the same realization. TOTALENV1 performs no quantifier elimination. FACTOR1 still gives one exact-decider layer for a fixed answer vector. Generic quantifier commutation remains \textbf{FALSIFIED/RETIRED}. Generic canonical binding is not free.
RESPSEL1 strengthens the hostile control: under the counterfactual P=NP, total polynomially balanced P-verifiable stage-response relations admit polynomial-time selectors. This does not supply the required source-matched evaluator unconditionally; it reinforces existence versus acquisition versus verification.
\section*{Full charge and TransportCert}
The time-level ledger remains
\[\tau_Q+\tau_{\rm norm}+\tau_{\rm census}+\tau_{\rm amp}+\tau_{\rm idx}+\tau_{\rm common\,compile}+\min\{\tau_D,\tau_R,\tau_F\}+\tau_{\rm write}+\tau_{\rm transport}.\]
TOTALENV1 only converts this constant-length sum into its polynomial-exponent maximum. TransportCert remains accounting only, not a lower-bound theorem.
\section*{PQ and cross-domain classification}
Construction-specific collision extraction or special soundness remains \textbf{ANALOGY/CANDIDATE} until an unconditional source-matched theorem is proved. Generic cryptographic binding and random-oracle arguments are \textbf{FAILED TRANSFER} for an unconditional separation. QBF proof-system Skolem extraction is a \textbf{KNOWN THEOREM} in its proof model but only an \textbf{ANALOGY/CANDIDATE} for arbitrary SAT computation because acquiring the proof representation is not free.
\section*{Barrier audit}
TOTALENV1 relativizes and survives algebraic-oracle augmentation. A semantic oracle can supply route outputs or common-charge objects. The theorem gives no necessity result forcing an arbitrary SAT decider to materialize Q, a selector, an evaluator, a proof, or a transcript. SAME\_REALIZATION acquisition, TC4 arbitrary-algorithm escape, nonrelativization and nonalgebrization remain \textbf{OPEN}.
\section*{WHAT CHANGED RELATIVE TO ROOT MAIN}
Root main.tex already records ENVELOPE1, RAMDOM1 and CROSSDOM1 and the additive full-charge ledger. CERTREG1 is verified but not yet reflected in the owned region. TOTALENV1 adds the exact full-exponent envelope \(\max\{G,E(K)\}\) and the bottleneck guard: improving route accounting cannot improve the total exponent while a common charge dominates.
\section*{ADOPTED / REJECTED / RETIRED PRIOR COMPONENTS}
\textbf{Adopted:} FACTOR1, QLEN1, RAMCOMP1, ENVELOPE1, RAMDOM1, CROSSDOM1, CERTREG1, EXPINV1, explicit evaluator acquisition, RESPSEL1 as a hostile closure, and separate common/route-specific charging.
\textbf{Rejected as free:} generic canonical binding, evaluator acquisition, proof acquisition, qualitative PH collapse as a numerical compiler, optimization over looser exponent certificates, and moving route-specific costs into the common ledger.
\textbf{FALSIFIED/RETIRED:} generic quantifier commutation and SGLAUD1 constant-query amplification remain retired.
\section*{STALE OR OVERSTATED CLAIMS FOUND}
No verified corpus artifact establishes an unrestricted SAT lower bound from the reconstruction envelope. Any reading of a smaller route exponent as a total-runtime improvement without comparing against common-charge exponent G is overstated and is corrected by TOTALENV1.
\section*{Next target}
Source-transcribe the actual vMS/SGL SAME\_REALIZATION tuple \((a,q,b,\ell,c_0,f,c_F,K)\) together with every common exponent. Only then evaluate \(\max\{G,E(K)\}\). If G dominates, stop optimizing route crossovers and attack the dominating common obligation instead.
\section*{References}
E. Viola, On Probabilistic Time versus Alternating Time, ECCC TR05-137, 2005.
D. van Melkebeek and N. Mocelin Sdroievski, Instance-Wise Hardness and Refutation versus Derandomization for Arthur-Merlin Protocols, Computational Complexity 34:15, 2025.
S. Akshay, O. Beyersdorff, S. Chakraborty, L. Kasche, M. Mahajan, and L. N. Spachmann, Proof Systems for QBF Synthesis: Extracting Skolem and Herbrand Functions, SAT 2026, LIPIcs 377, Article 3.
\end{document}
% END_OF_VERIFIED_REPORT::LSPV-3::TOTALENV1